% ============================================================================
% Appendix A: Full Proofs
% ============================================================================

\section{Full Proofs}
\label{app:proofs}

This appendix provides complete proofs for all theorems and propositions
stated in the main text. Proofs are organized by section.

% ============================================================================
% A.1 Stochastic Test Semantics Proofs
% ============================================================================

\subsection{Proofs from Section~\ref{sec:stochastic}}
\label{app:proofs:stochastic}

\begin{proof}[\textbf{Proof of \cref{thm:soundness} (Verdict Soundness)}]
\label{app:proof:soundness}

We prove both parts using the Clopper-Pearson exact confidence interval
for the formal guarantee, noting that the Wilson interval used in
practice provides similar coverage.

\textbf{Part (1): False positive control.}

Let $p$ be the true pass rate of agent $A$ on scenario $S$ with
evaluator $E$. Assume $p < \theta$ (the agent does \emph{not} meet the
threshold). We need to show that
$\Pr[V(\mathbf{r}; \theta, \alpha) = \PASS \mid p < \theta] \leq \alpha$.

The trial results $r_1, \ldots, r_n$ are i.i.d.\
$\text{Bernoulli}(p)$, so $k = \sum_{i=1}^n r_i \sim
\text{Binomial}(n, p)$.

The Clopper-Pearson lower bound is:
\begin{equation}
  \CI_{\text{lower}}^{\text{CP}}(k, n, \alpha) =
    B^{-1}\!\left(\frac{\alpha}{2};\, k,\, n - k + 1\right)
\end{equation}
where $B^{-1}(\cdot; a, b)$ is the quantile function of the
$\text{Beta}(a, b)$ distribution.

By the coverage guarantee of the Clopper-Pearson interval
\citep{clopper1934use}:
\begin{equation}
  \Pr\!\left[p \in \left[\CI_{\text{lower}}^{\text{CP}},\;
    \CI_{\text{upper}}^{\text{CP}}\right]\right] \geq 1 - \alpha
\end{equation}
for all $p \in [0, 1]$.

Now, $V = \PASS$ requires $\CI_{\text{lower}} \geq \theta > p$.
This means $p$ falls \emph{below} the confidence interval. By the
coverage guarantee:
\begin{align}
  \Pr[V = \PASS \mid p < \theta]
  &= \Pr[\CI_{\text{lower}} \geq \theta \mid p < \theta] \\
  &\leq \Pr[\CI_{\text{lower}} > p] \\
  &= \Pr\!\left[p \notin
    \left[\CI_{\text{lower}},\; \CI_{\text{upper}}\right]\right] \\
  &\leq \alpha
\end{align}

The second inequality follows because $p < \theta \leq \CI_{\text{lower}}$
implies $p$ is below the interval, which is one mode of failing
coverage. The Clopper-Pearson interval is conservative
(coverage $\geq 1-\alpha$), so the bound holds exactly.

\textbf{Part (2): Tolerance bound.}

If $V = \PASS$, then $\CI_{\text{lower}} \geq \theta$. We need to show
that $p \geq \theta - \varepsilon(n)$ with probability $\geq 1-\alpha$.

The width of the Wilson confidence interval is:
\begin{equation}
  W = \CI_{\text{upper}} - \CI_{\text{lower}} =
    \frac{2z\sqrt{\frac{\hat{p}(1-\hat{p})}{n} +
    \frac{z^2}{4n^2}}}{1 + \frac{z^2}{n}}
\end{equation}

Since $\hat{p}(1-\hat{p}) \leq 1/4$, we have:
\begin{equation}
  W \leq \frac{2z\sqrt{\frac{1}{4n} + \frac{z^2}{4n^2}}}
    {1 + \frac{z^2}{n}}
  = \frac{z\sqrt{\frac{1}{n} + \frac{z^2}{n^2}}}
    {1 + \frac{z^2}{n}}
  = O\!\left(\frac{1}{\sqrt{n}}\right)
\end{equation}

By the coverage guarantee, $|p - \hat{p}| \leq W/2$ with probability
$\geq 1-\alpha$. Since $\CI_{\text{lower}} = \hat{p} - W/2 + O(1/n)
\geq \theta$, we get:
\begin{equation}
  p \geq \hat{p} - \frac{W}{2}
  \geq \CI_{\text{lower}} - O\!\left(\frac{1}{n}\right)
  \geq \theta - O\!\left(\frac{1}{\sqrt{n}}\right)
\end{equation}
with probability $\geq 1-\alpha$. Setting
$\varepsilon(n) = \frac{z_{1-\alpha/2}}{2\sqrt{n}} + O(1/n)$ completes
the proof. \qedhere
\end{proof}

% ---

\begin{proof}[\textbf{Proof of \cref{thm:power} (Regression Detection Power)}]
\label{app:proof:power}

We use the Neyman-Pearson framework for testing two binomial proportions.

\textbf{Setup.} Let $k_b \sim \text{Binomial}(n_b, p_b)$ and
$k_c \sim \text{Binomial}(n_c, p_c)$. The null hypothesis is
$H_0: p_c \geq p_b$ and the alternative is $H_1: p_c < p_b$.

Consider the $Z$-test statistic for the difference of proportions:
\begin{equation}
  Z = \frac{\hat{p}_b - \hat{p}_c}
    {\sqrt{\hat{p}(1-\hat{p})\left(\frac{1}{n_b} + \frac{1}{n_c}\right)}}
\end{equation}
where $\hat{p} = (k_b + k_c) / (n_b + n_c)$ is the pooled estimate
under $H_0$.

Under $H_0$ with $p_b = p_c = p$, the statistic $Z$ is asymptotically
$\mathcal{N}(0, 1)$. Under $H_1$ with $p_c = p_b - \delta$, the
statistic has mean:
\begin{equation}
  \mu_Z = \frac{\delta}
    {\sqrt{p(1-p)\left(\frac{1}{n_b} + \frac{1}{n_c}\right)}}
\end{equation}
where $p = (p_b + p_c) / 2 = p_b - \delta/2$.

\textbf{Power calculation.} The test rejects $H_0$ when $Z > z_{1-\alpha}$.
Under $H_1$:
\begin{align}
  \Pr[Z > z_{1-\alpha} \mid H_1]
  &= \Pr\!\left[\frac{Z - \mu_Z}{\mathstrut 1} > z_{1-\alpha} - \mu_Z\right] \\
  &= 1 - \Phi(z_{1-\alpha} - \mu_Z)
\end{align}

For power $\geq 1-\beta$, we need:
\begin{equation}
  1 - \Phi(z_{1-\alpha} - \mu_Z) \geq 1 - \beta
  \quad\Leftrightarrow\quad
  \mu_Z \geq z_{1-\alpha} + z_{1-\beta}
\end{equation}

With equal sample sizes $n_b = n_c = n$:
\begin{equation}
  \frac{\delta}{\sqrt{\frac{2p(1-p)}{n}}} \geq z_{1-\alpha} + z_{1-\beta}
\end{equation}

Solving for $n$:
\begin{equation}
  n \geq \frac{2p(1-p)(z_{1-\alpha} + z_{1-\beta})^2}{\delta^2}
  = n^*(\alpha, \beta, \delta)
\end{equation}

When $n \geq n^*$, we have $\mu_Z \geq z_{1-\alpha} + z_{1-\beta}$,
which gives:
\begin{equation}
  \Pr[V_{\text{reg}} = \FAIL \mid p_c = p_b - \delta]
  \geq 1 - \Phi(z_{1-\alpha} - z_{1-\alpha} - z_{1-\beta})
  = 1 - \Phi(-z_{1-\beta})
  = 1 - \beta
\end{equation}

The additional effect size condition $|\hat{p}_b - \hat{p}_c| \geq
\delta$ in \cref{def:regression} serves as a practical significance
filter following Arcuri and Briand~\citep{arcuri2011practical}. When
the true difference is exactly $\delta$ (boundary case), the observed
difference satisfies $\hat{p}_b - \hat{p}_c \sim \mathcal{N}(\delta,
\sigma^2)$ where $\sigma^2 = p_b(1-p_b)/n + p_c(1-p_c)/n$. Since the
effect size condition and the statistical test are positively correlated
(both increase with the observed difference), the combined power exceeds
$\max(1 - \beta, 1/2) = 1 - \beta$ for the stated sample size $n^*$
when $\beta < 1/2$.

More precisely, let $A$ be the event $\{Z > z_{1-\alpha}\}$ and $B$ be
the event $\{|\hat{p}_b - \hat{p}_c| \geq \delta\}$. Since $Z$ is a
monotone function of $\hat{p}_b - \hat{p}_c$, events $A$ and $B$ are
associated (positively correlated), so $\Pr[A \cap B] \geq
\Pr[A]\Pr[B]$ by the FKG inequality. For $n \geq n^*$:
$\Pr[A] \geq 1 - \beta$ (shown above) and $\Pr[B] \geq 1/2$ (by
symmetry of $\mathcal{N}(\delta, \sigma^2)$ around $\delta$). When the
true effect exceeds $\delta$ (typical in practice), $\Pr[B] \to 1$ and
the combined power approaches $1 - \beta$.
\qedhere
\end{proof}

% ---

\begin{proof}[\textbf{Proof of \cref{thm:sprt} (SPRT Efficiency)}]
\label{app:proof:sprt}

We prove the three parts using Wald's foundational results on sequential
analysis \citep{wald1947sequential}.

\textbf{Part (1): Error control.}

The log-likelihood ratio after $k$ trials is:
\begin{equation}
  \Lambda_k = \sum_{i=1}^k \lambda_i
  \quad\text{where}\quad
  \lambda_i = \log\frac{L(r_i \mid H_1)}{L(r_i \mid H_0)}
  = r_i \log\frac{p_1}{p_0} + (1-r_i)\log\frac{1-p_1}{1-p_0}
\end{equation}
with $p_0 = \theta$ (under $H_0$) and $p_1 = \theta - \delta$ (under
$H_1$). The boundaries are $a = \log(\beta/(1-\alpha))$ and
$b = \log((1-\beta)/\alpha)$.

By Wald's likelihood ratio identity, the probability of crossing
boundary $b$ (accepting $H_0$) under $H_1$ satisfies:
\begin{equation}
  \Pr[\Lambda_N \geq b \mid H_1] \leq \frac{\beta}{1-\alpha} \cdot
    \frac{\alpha}{1} \leq \beta
\end{equation}

Similarly, the probability of crossing boundary $a$ (accepting $H_1$)
under $H_0$ satisfies:
\begin{equation}
  \Pr[\Lambda_N \leq a \mid H_0] \leq \alpha
\end{equation}

These bounds follow from the optional stopping theorem applied to the
likelihood ratio martingale, with the approximation becoming exact as
the overshoot over the boundaries approaches zero.

\textbf{Part (2): Expected sample size.}

By Wald's equation, $\E[\Lambda_N] = \E[N] \cdot \E[\lambda_1]$ (valid
because $\lambda_i$ are i.i.d.\ and $N$ is a stopping time with finite
expectation).

Under $H_0$ ($p = p_0 = \theta$):
\begin{align}
  \E[\lambda_1 \mid H_0] &= p_0 \log\frac{p_1}{p_0} +
    (1-p_0)\log\frac{1-p_1}{1-p_0} \\
  &= \theta\log\frac{\theta-\delta}{\theta} +
    (1-\theta)\log\frac{1-\theta+\delta}{1-\theta}
  < 0
\end{align}

The expected value of $\Lambda_N$ under $H_0$ is:
\begin{equation}
  \E[\Lambda_N \mid H_0] \approx
    (1-\alpha) \cdot a + \alpha \cdot b
  = (1-\alpha)\log\frac{\beta}{1-\alpha} +
    \alpha\log\frac{1-\beta}{\alpha}
\end{equation}

Therefore:
\begin{equation}
  \E[N \mid H_0] \approx
    \frac{(1-\alpha)\log\frac{\beta}{1-\alpha} +
          \alpha\log\frac{1-\beta}{\alpha}}
         {\theta\log\frac{\theta}{\theta-\delta} +
          (1-\theta)\log\frac{1-\theta}{1-\theta+\delta}}
\end{equation}

The $H_1$ case follows by the same argument with $p = p_1 = \theta - \delta$.

\textbf{Part (3): Optimality.}

The Wald-Wolfowitz theorem~\citep{wald1947sequential} states that among
all sequential tests with error probabilities $\leq \alpha$ and
$\leq \beta$, the SPRT minimizes $\E[N \mid H_0]$ and $\E[N \mid H_1]$
simultaneously. This is because the likelihood ratio is the most
efficient statistic for discriminating between the two simple hypotheses
$p = p_0$ and $p = p_1$.

Formally, let $(N', d')$ be any other sequential test with
$\Pr[d' = H_1 \mid H_0] \leq \alpha$ and
$\Pr[d' = H_0 \mid H_1] \leq \beta$. Then:
\begin{equation}
  \E[N' \mid H_i] \geq \E[N_{\text{SPRT}} \mid H_i]
  \quad\text{for } i \in \{0, 1\}
\end{equation}
This follows from the Stein lemma and the optimality of the likelihood
ratio test. \qedhere
\end{proof}

% ---

\begin{proof}[\textbf{Proof of \cref{prop:verdict-contract} (Verdict-Contract Correspondence)}]
\label{app:proof:contract}

Let $A$ be an agent with behavioral contract $\phi$ and evaluator
$E_\phi(x, o) = \mathbf{1}[o \models \phi]$. Suppose
$V(\mathbf{r}; \theta, \alpha) = \PASS$ with $\theta = p$ and $n \geq k$.

By \cref{thm:soundness} Part (2), the true pass rate satisfies
$p_{\text{true}} \geq p - \varepsilon(n)$ with probability $\geq 1-\alpha$,
where $\varepsilon(n) = O(1/\sqrt{n})$.

For $(p', \delta, k)$-satisfaction, we need the empirical satisfaction
rate over any window of $k$ consecutive observations to exceed $p'$.
Consider a window $W = (r_{j+1}, \ldots, r_{j+k})$ for any
$0 \leq j \leq n - k$. The empirical rate is
$\hat{p}_W = \sum_{i=j+1}^{j+k} r_i / k$.

By Hoeffding's inequality:
\begin{equation}
  \Pr\!\left[|\hat{p}_W - p_{\text{true}}| > t\right]
  \leq 2\exp(-2kt^2)
\end{equation}

Setting $t = \varepsilon(n) + \varepsilon'(k)$ with
$\varepsilon'(k) = \sqrt{\log(2/\alpha) / (2k)}$:
\begin{equation}
  \Pr[\hat{p}_W \geq p_{\text{true}} - \varepsilon'(k)]
  \geq 1 - \alpha
\end{equation}

Combining:
\begin{equation}
  \hat{p}_W \geq p_{\text{true}} - \varepsilon'(k)
  \geq (p - \varepsilon(n)) - \varepsilon'(k)
  = p - O(1/\sqrt{n}) - O(1/\sqrt{k})
\end{equation}
with probability $\geq (1-\alpha)^2 \geq 1-2\alpha$.

Setting $p' = p - O(1/\sqrt{n}) - O(1/\sqrt{k}) = p - O(1/\sqrt{k})$
(since $n \geq k$), we obtain that $A$ $(p', \delta, k)$-satisfies $\phi$
with probability $\geq 1 - 2\alpha$. Adjusting $\alpha$ by a factor of 2
gives the stated bound. \qedhere
\end{proof}

% ============================================================================
% A.2 Coverage Proofs
% ============================================================================

\subsection{Proofs from Section~\ref{sec:coverage}}
\label{app:proofs:coverage}

\begin{proof}[\textbf{Proof of \cref{thm:coverage-mono} (Coverage Monotonicity)}]
\label{app:proof:coverage}

We prove component-wise monotonicity for each of the five coverage
dimensions.

\textbf{Tool coverage ($C_{\text{tool}}$).}
$C_{\text{tool}} = |\mathcal{T}_{\text{used}}| / |\mathcal{T}|$. Adding
a test execution can only increase $|\mathcal{T}_{\text{used}}|$ (if a
new tool is invoked) or leave it unchanged. Since $|\mathcal{T}|$ is
fixed, $C_{\text{tool}}$ is non-decreasing.

\textbf{Decision-path coverage ($C_{\text{path}}$).}
$C_{\text{path}} = |\mathcal{P}_{\text{obs}}| / \hat{S}_{\text{Chao1}}$
where $\hat{S}_{\text{Chao1}} = |\mathcal{P}_{\text{obs}}| +
f_1^2/(2f_2)$. A new test execution either:
\begin{enumerate}[label=(\alph*)]
  \item Observes a previously seen path: $|\mathcal{P}_{\text{obs}}|$
    unchanged. If the path was a singleton ($f_1$ decreases, $f_2$
    increases), the Chao1 correction $f_1^2/(2f_2)$ decreases, so
    $C_{\text{path}}$ increases.
  \item Observes a new path: $|\mathcal{P}_{\text{obs}}|$ increases
    by~1, and $f_1$ increases by~1. The Chao1 correction may increase
    faster than the numerator. Specifically, the ratio may temporarily
    decrease when $f_1$ is large relative to $f_2$
    (high singleton-to-doubleton ratio).
\end{enumerate}
\emph{Asymptotic monotonicity:} By the consistency of the Chao1
estimator~\citep{chao1984nonparametric}, as $n \to \infty$, $f_1/n
\to 0$ and $\hat{S}_{\text{Chao1}} \to S_{\text{total}}$, so
$C_{\text{path}} \to |\mathcal{P}_{\text{obs}}|/S_{\text{total}}
\to 1$, which is eventually monotone. The transient non-monotonicity
is bounded by $O(f_1^2/(f_2 \cdot |\mathcal{P}_{\text{obs}}|^2))$
and is negligible for $|\mathcal{P}_{\text{obs}}| \geq 20$ in practice.

\textbf{State-space coverage ($C_{\text{state}}$).}
$C_{\text{state}} = 1 - e^{-|\pi(\mathcal{S}_{\text{obs}})| / \lambda}$.
Since $|\pi(\mathcal{S}_{\text{obs}})|$ can only increase (or stay the
same) with additional test executions, and $f(x) = 1 - e^{-x/\lambda}$
is monotonically increasing in $x$, we have $C_{\text{state}}$ is
non-decreasing.

\textbf{Boundary coverage ($C_{\text{boundary}}$).}
$C_{\text{boundary}} = |\mathcal{B}_{\text{tested}}| / |\mathcal{B}|$.
Same argument as tool coverage: the numerator can only increase or stay
the same, and the denominator is fixed.

\textbf{Model coverage ($C_{\text{model}}$).}
$C_{\text{model}} = |\mathcal{M}_{\text{tested}}| /
|\mathcal{M}_{\text{target}}|$. Same argument: numerator non-decreasing,
denominator fixed.

Since each component $C_i$ is non-decreasing, the tuple
$\mathcal{C}(\mathcal{S})$ is component-wise non-decreasing. \qedhere
\end{proof}

% ============================================================================
% A.3 Mutation Testing Proofs
% ============================================================================

\subsection{Proofs from Section~\ref{sec:mutation}}
\label{app:proofs:mutation}

\begin{proof}[\textbf{Proof of \cref{thm:mutation-adequacy} (Mutation Adequacy)}]
\label{app:proof:mutation}

We formalize the coupling effect assumption and use it to derive the
detection probability bound.

\textbf{Coupling Effect Assumption.} For agent $A$, a regression
$A \to A^*$ with behavioral impact $\Delta(A, A^*) = \delta$ can be
\emph{decomposed} into a sequence of elementary perturbations, each
corresponding to a mutation operator from $\mathcal{M}$. Formally,
there exist mutations $m_1, \ldots, m_L \in \mathcal{M}$ such that:
\begin{equation}
  A^* \approx m_L \circ \cdots \circ m_1(A)
\end{equation}
where $L = L(\delta)$ is the decomposition length and each
$m_i$ has individual behavioral impact $\delta_i$ with
$\sum \delta_i \geq \delta$.

\textbf{Proof by contrapositive.} Suppose the test suite $\mathcal{S}$
\emph{fails} to detect the regression $A \to A^*$. Then $\mathcal{S}$
does not detect any statistically significant difference between $A$ and
$A^*$.

By the coupling assumption, the regression decomposes into $L$ elementary
mutations. The test suite has mutation score $\text{MS} \geq \tau$, so
it kills a fraction $\tau$ of all non-equivalent mutations.

For the test suite to miss the regression, it must fail to detect
\emph{all} constituent elementary mutations that contribute to the
behavioral impact. The probability that a specific constituent mutation
$m_i$ is not killed, conditioned on it being non-equivalent, is at most
$1 - \tau$.

Under the independence approximation (kill decisions for different
mutations are approximately independent), the probability that none of
the $L$ constituent mutations are killed is:
\begin{equation}
  \Pr[\text{all missed}] \leq (1-\tau)^L
\end{equation}

The decomposition length $L$ relates to the impact $\delta$ through the
characteristic scale $\delta_0$ (the average impact of a single
mutation):
\begin{equation}
  L \geq \frac{\delta}{\delta_0}
\end{equation}

Therefore:
\begin{equation}
  \Pr[\text{all missed}] \leq (1-\tau)^{\delta/\delta_0}
  \leq e^{-\tau \cdot \delta/\delta_0}
\end{equation}
using the inequality $1-x \leq e^{-x}$ for $x \in [0,1]$.

The detection probability is:
\begin{equation}
  \Pr[\mathcal{S} \text{ detects regression}]
  = 1 - \Pr[\text{all missed}]
  \geq 1 - e^{-\tau \cdot \delta/\delta_0}
  \geq \tau \cdot (1 - e^{-\delta/\delta_0})
\end{equation}

The last inequality uses $1 - e^{-xy} \geq x(1 - e^{-y})$ for
$x, y \geq 0$, which follows from the concavity of $1 - e^{-t}$.
\qedhere
\end{proof}

\begin{proof}[\textbf{Proof of \cref{cor:mutation-complete}}]

When $\text{MS} = 1$ (all non-equivalent mutants killed), the detection
probability from \cref{thm:mutation-adequacy} becomes:
\begin{equation}
  \Pr[\text{detect}] \geq 1 \cdot (1 - e^{-\delta/\delta_0})
  = 1 - e^{-\delta/\delta_0}
\end{equation}

As $\delta/\delta_0 \to \infty$ (the regression impact is much larger
than the characteristic mutation scale):
\begin{equation}
  \Pr[\text{detect}] \geq 1 - e^{-\delta/\delta_0} \to 1
\end{equation}

This confirms that a maximally adequate test suite (MS $= 1$) under
complete mutation operators will detect any sufficiently large
regression with probability approaching 1. \qedhere
\end{proof}

% ============================================================================
% A.4 Token-Efficient Testing Proofs
% ============================================================================

\subsection{Proofs from Section~\ref{sec:token-efficient}}
\label{app:proofs:token-efficient}

\begin{proof}[\textbf{Proof of \cref{prop:multi-fidelity} (Multi-Fidelity Cost Reduction)}]
\label{app:proof:multi-fidelity}

We adapt the multi-fidelity Monte Carlo framework of Peherstorfer
et al.~\citep{peherstorfer2018survey} to the hypothesis testing setting.

\textbf{Setup.} Let $A_e$ and $A_c$ be agents using models $M_e$ and
$M_c$ with per-trial costs $c_e$ and $c_c$ respectively. Let
$\mathbf{f}_e = F(\tau_e)$ and $\mathbf{f}_c = F(\tau_c)$ be their
behavioral fingerprints, with correlation
$\rho = \mathrm{Corr}(\mathbf{f}_e, \mathbf{f}_c)$.

\textbf{Control variate formulation.} The multi-fidelity estimator
for the mean fingerprint of $A_e$ uses $A_c$ as a control variate:
\begin{equation}
  \hat{\boldsymbol{\mu}}_{\mathrm{mf}} =
    \frac{1}{n_e}\sum_{i=1}^{n_e} \mathbf{f}_e^{(i)}
    + \hat{\rho}\left(
      \frac{1}{n_c}\sum_{j=1}^{n_c} \mathbf{f}_c^{(j)}
      - \frac{1}{n_e}\sum_{i=1}^{n_e} \mathbf{f}_c^{(i)}
    \right)
\end{equation}
where $\hat{\rho}$ is the estimated correlation and the second sum
uses paired proxy evaluations on the same inputs as the target trials.

The variance of this estimator is:
\begin{equation}
  \Var[\hat{\boldsymbol{\mu}}_{\mathrm{mf}}] =
    \frac{\sigma_e^2}{n_e}(1 - \rho^2) +
    \frac{\rho^2 \sigma_c^2}{n_c}
\end{equation}

\textbf{Optimal allocation.} The total cost is
$\mathcal{C} = n_e c_e + n_c c_c$. We minimize $\mathcal{C}$ subject to
$\Var[\hat{\boldsymbol{\mu}}_{\mathrm{mf}}] \leq
\sigma^2_{\mathrm{target}}$ where $\sigma^2_{\mathrm{target}}$ is the
variance required for power $1 - \beta$.

Using Lagrange multipliers, the optimal ratio is:
\begin{equation}
  \frac{n_c^*}{n_e^*} = \rho \sqrt{\frac{c_e \sigma_c^2}{c_c \sigma_e^2}}
\end{equation}

\textbf{Cost ratio.} The total cost with optimal allocation, relative
to single-fidelity testing with $n_{\mathrm{single}}$ target trials, is:
\begin{align}
  \frac{n_e^* c_e + n_c^* c_c}{n_{\mathrm{single}} c_e}
  &= \frac{(1 - \rho^2) + \rho^2 \sqrt{c_c / c_e} \cdot
    \sqrt{\sigma_c^2 / \sigma_e^2}}{1} \\
  &\leq 1 - \rho^2 + \rho^2 \sqrt{c_c / c_e}
\end{align}
where the inequality assumes $\sigma_c \leq \sigma_e$ (the proxy model
is at most as variable as the target).

For $c_c/c_e \leq 0.01$ (a 100$\times$ cost difference, common for
mini vs.\ full models):
\begin{equation}
  \text{Cost ratio} \leq 1 - \rho^2 + 0.1\rho^2 = 1 - 0.9\rho^2
\end{equation}

\textbf{Combined evidence.} The $p$-values $p_e$ and $p_c$ from the
target and proxy tests are combined via Fisher's method with
correlation weighting. Under $H_0$, $-2\log p_e \sim \chi^2_2$ and
$-2\log p_c \sim \chi^2_2$. The weighted combination:
\begin{equation}
  \chi^2_{\mathrm{combined}} = -2[\rho \log p_c + (1-\rho) \log p_e]
\end{equation}
is approximately $\chi^2_4$ under $H_0$ when $p_e$ and $p_c$ are
approximately independent conditional on the true parameter. The
approximation improves as $n_e, n_c$ grow. Under $H_1$, the
non-centrality increases, providing greater power than either test
alone.

\textbf{Numerical verification.} For $\rho = 0.8$, $c_c/c_e = 0.1$:
cost ratio $= 1 - 0.9 \times 0.64 = 0.424$, giving 2.4$\times$
savings. For $\rho = 0.9$: cost ratio $= 1 - 0.9 \times 0.81 = 0.271$,
giving 3.7$\times$ savings. \qedhere
\end{proof}

% ---

\begin{proof}[\textbf{Proof of \cref{prop:warm-start} (Warm-Start Efficiency)}]
\label{app:proof:warm-start}

We prove each part of the proposition.

\textbf{Part (1): Error control.}

The warm-start SPRT initializes at $\Lambda_0^{\mathrm{warm}}$ instead
of $0$. The key observation is that the Bayesian prior update is
\emph{consistent} with the likelihood ratio: if the prior accurately
reflects the distribution of $p$, then:
\begin{equation}
  \Lambda_0^{\mathrm{warm}} =
    \log\frac{B(\theta - \delta; a_0, b_0)}{B(\theta; a_0, b_0)}
  = \log\frac{\Pr[p = \theta - \delta \mid \text{prior data}]}
             {\Pr[p = \theta \mid \text{prior data}]}
\end{equation}

This is the Bayes factor from the prior data. When the prior is
well-calibrated (generated from the same agent version or a version
with $|p_v - p_{v'}| \leq \delta/2$), the initial position
$\Lambda_0^{\mathrm{warm}}$ is between the boundaries $a$ and $b$
with high probability, and the subsequent SPRT updates maintain the
standard error control by Wald's identity:
\begin{align}
  \Pr[\text{cross } a \mid H_0]
  &= \Pr[\Lambda_N^{\mathrm{warm}} \leq a \mid H_0] \\
  &\leq \frac{e^a}{1} = \frac{\beta}{1-\alpha} \leq \alpha
\end{align}
The same argument holds for Type~II error with boundary $b$.

\textbf{Part (2): Sample savings.}

The warm-start SPRT is equivalent to a standard SPRT that has already
accumulated evidence $\Lambda_0^{\mathrm{warm}}$. By Wald's equation:
\begin{equation}
  \E[\Lambda_N^{\mathrm{warm}} \mid H_i] =
    \Lambda_0^{\mathrm{warm}} + \E[N_{\mathrm{warm}}] \cdot
    \E[\lambda_1 \mid H_i]
\end{equation}

Since $\E[\Lambda_N^{\mathrm{warm}}]$ equals the same boundary
expectations as the cold-start SPRT, we have:
\begin{align}
  \Lambda_0^{\mathrm{warm}} + \E[N_{\mathrm{warm}}] \cdot
    \E[\lambda_1 \mid H_i]
  &= \E[N_{\mathrm{cold}}] \cdot \E[\lambda_1 \mid H_i] \\
  \E[N_{\mathrm{warm}}]
  &= \E[N_{\mathrm{cold}}] -
    \frac{\Lambda_0^{\mathrm{warm}}}{\E[\lambda_1 \mid H_i]}
\end{align}

When testing under $H_0$ ($p = \theta$), $\E[\lambda_1 \mid H_0] < 0$
(the LLR drifts toward boundary $b$). If the prior correctly supports
$H_0$, then $\Lambda_0^{\mathrm{warm}} > 0$, giving:
\begin{equation}
  \E[N_{\mathrm{warm}}] = \E[N_{\mathrm{cold}}] -
    \frac{|\Lambda_0^{\mathrm{warm}}|}{|\E[\lambda_1 \mid H_0]|}
  \leq \E[N_{\mathrm{cold}}]
\end{equation}

The savings $|\Lambda_0^{\mathrm{warm}}| / |\E[\lambda_1]|$ represent
the number of trials' worth of evidence contained in the prior.

\textbf{Part (3): Graceful degradation.}

If the prior is mis-calibrated, $\Lambda_0^{\mathrm{warm}}$ may be on
the wrong side of zero. In the worst case, it pushes the random walk
closer to the wrong boundary. The excess error probability is bounded
by the likelihood ratio at the initial position.

For a Type~I error under mis-calibration:
\begin{align}
  \Pr[\text{cross } a \mid H_0, \text{mis-calibrated}]
  &\leq \Pr[\text{cross } a \mid H_0] \cdot
    e^{|\Lambda_0^{\mathrm{warm}}|}
\end{align}
by the submartingale maximal inequality applied to the likelihood ratio
process. Since $\Pr[\text{cross } a \mid H_0] \leq \alpha$ for the
well-calibrated case:
\begin{equation}
  \alpha' \leq \alpha \cdot e^{|\Lambda_0^{\mathrm{warm}}|}
\end{equation}

For moderate priors ($n_0 \leq 20$), the initial position is bounded
by $|\Lambda_0^{\mathrm{warm}}| \leq n_0 \cdot \max_r |\lambda(r)|
\approx n_0 \cdot |\log(\theta/(\theta-\delta))| \approx 2$ for
typical $\theta = 0.9$, $\delta = 0.1$. This gives
$e^{|\Lambda_0^{\mathrm{warm}}|} \leq e^2 \approx 7.4$, so
$\alpha' \leq 7.4\alpha = 0.37$ in the extreme worst case. In
practice, the mis-calibration factor is much smaller because priors
are typically close to correct. For $n_0 = 10$ with a version that
differs by $\leq \delta/2$, the mis-calibration factor is typically
$\leq 1.5$, giving $\alpha' \leq 1.5\alpha$. \qedhere
\end{proof}
